\section{Roadmap Update after Experiment 11A}
\label{sec:roadmap-post-11a}

Experiment~11A changes the status of the descent-aware branch substantially. The question is no longer whether principled event sizing can matter: the exact global descent oracle strongly outperforms the best heuristic/global schedule on the controlled diagonal quadratic. The unresolved question is whether enough information for a global descent certificate can be obtained in a communication-efficient federated system.

The updated mechanism status is:
\begin{enumerate}
  \item \textbf{Fixed/global scheduled jump resolution:} remains the strongest currently deployable reference. The schedule from Experiment~10C should therefore remain the practical baseline in subsequent robustness experiments.
  \item \textbf{Global descent-aware jump sizing:} strong oracle success. This is now the principal upper bound for what sparse event sizing could achieve.
  \item \textbf{Pure local descent sizing:} rejected as a general federated rule. Even exact client-local gradients produce globally harmful updates under modest heterogeneity.
  \item \textbf{Heterogeneity-aware descent certificate:} strong oracle success. A valid bound on local--global gradient disagreement is sufficient in the current benchmark, but estimating that bound without dense communication is open.
  \item \textbf{Raw first-passage/event-derived sizing:} not sufficient. Event timing estimates local gradient magnitude surprisingly well on average, but sign/magnitude errors and client heterogeneity prevent a reliable global rule, especially for slow clients.
  \item \textbf{Event-derived sizing with valid uncertainty/disagreement margins:} oracle success. This identifies uncertainty quantification, not the timing formula itself, as the missing ingredient.
\end{enumerate}

\subsection{Chronology after Experiment 11A}

The numbering convention described in \cref{sec:numbering-chronology-note} remains in force. The next executed study should return to \textbf{Experiment~10D}, because the current low-bit coordinate encoder still needs the planned non-axis-aligned geometry stress test. Experiment~10D should use
\begin{equation}
  \boxed{\text{normalized full-reset coordinate LIF + tuned global }q(t)}
\end{equation}
as the deployable reference and should carry the exact global descent rule as an oracle upper bound where useful. The unimplemented descent certificate should not be promoted into the core optimizer before its information requirements are resolved.

After the geometry gate, the next descent-development stage should be denoted \textbf{Experiment~11B}. Its target is no longer generic ``adaptive $q$'' but the specific information problem exposed by 11A:
\begin{equation}
  \boxed{\text{estimate or maintain a sparse conservative bound on global event alignment and uncertainty.}}
\end{equation}
Candidate directions include client-specific or coordinate-specific heterogeneity envelopes, server-side estimates based on historical accepted events, confidence bounds derived from event timing plus model-drift information, clustered/client-similarity bounds, or occasional low-rate calibration messages.

\subsection{A second opportunity: descent-aware event generation}

Experiment~11A also shows that event sizing alone is not the final architecture. The global oracle accepts only about $1.5\%$ of transmitted threshold crossings in the primary benchmark. Most candidate events are generated by persistent heterogeneous local gradients after the global model is already close to optimal. If a conservative certificate could be evaluated before transmission, the same descent information might be used to modify the trigger itself,
\begin{equation}
  |z_{i,j}|\geq\Delta_j
  \quad\text{and}\quad
  \underline a_{i,j}>0,
\end{equation}
rather than transmitting every threshold crossing and rejecting most of them at the server.

This possible joint trigger-and-sizing mechanism is deliberately postponed until the certificate problem is understood. Otherwise a communication reduction could be obtained simply by silencing uncertain clients or coordinates, repeating the fairness failure observed in Experiment~10B. Any future descent-aware gate must therefore report acceptance/transmission rates by client compute period and by coordinate scale.

\subsection{Updated paper-level target}

The strongest prospective contribution is now more specific than the original ``LIF optimizer'' framing:
\begin{quote}
A neuromorphic federated optimizer in which local gradients are temporally encoded as sparse threshold events, while event acceptance and amplitude are determined from a conservative descent certificate constructed from event-derived evidence and federated uncertainty information.
\end{quote}
The current experiments support the communication encoder and the potential value of descent-aware sizing separately. They do not yet establish a practical method for constructing the global certificate. That gap should remain explicit in all future interpretations.
